% ensemble_readout_optimization.tex
% Optimal Ensemble Size and Qutrit Readout Protocol for NV-Based Qudit Computing
% NVQbit Theory Module — 2026-03-19
%
% Compilation: pdflatex + bibtex + pdflatex + pdflatex
% References: ../training/references.bib + local ensemble_readout.bib

\documentclass[11pt,letterpaper]{article}

\usepackage[margin=1in]{geometry}
\usepackage{amsmath,amssymb,amsfonts}
\usepackage{physics}
\usepackage{graphicx}
\usepackage{booktabs}
\usepackage{hyperref}
\usepackage{cleveref}
\usepackage[numbers,sort&compress]{natbib}
\usepackage{xcolor}

\newcommand{\needref}[1]{\textcolor{red}{\textbf{[NEEDREF: #1]}}}

\title{Optimal Ensemble Size and Qutrit Readout Protocol \\
for NV-Based Qudit Architecture}
\author{NVQbit Theory Module}
\date{2026-03-19}

\begin{document}
\maketitle

\begin{abstract}
We derive the minimum ensemble size $N$ of NV centers per logical qudit
required to achieve target readout fidelity for a ternary ($d=3$) quantum
computation scheme.  Two readout protocols are analyzed: (i) direct
three-level discrimination exploiting the differential fluorescence contrast
$C_2$ between $m_s = -1$ and $m_s = +1$, and (ii) sequential binary readout
using only the standard $m_s = 0$ vs.\ $m_s = \pm 1$ contrast $C_1 \approx 0.3$.
We show that sequential binary readout is strictly superior, eliminating
dependence on the vanishingly small $C_2$ and reducing $N_{\min}$ by orders of
magnitude.  We derive the optimal ensemble size $k^*$ that maximizes a
computational utility metric (accessible Hilbert-space dimension weighted by
readout fidelity), and establish practical upper bounds on the qudit base $d$
for three photon-collection technology tiers.
\end{abstract}

\tableofcontents

%% ==================================================================
\section{State-of-the-Art NV Photon Parameters}
\label{sec:params}
%% ==================================================================

All parameters below are drawn from published literature; none are from the
NVQbit project's own toy-setup measurements ($\eta = 14~\mu\text{T}/\!\sqrt{\text{Hz}}$).

\subsection{Photon Count Rate $R$}

\begin{table}[h]
\centering
\caption{Photon collection rates from single NV centers.}
\label{tab:rates}
\begin{tabular}{lccl}
\toprule
\textbf{Collection method} & \textbf{$R$ (cps)} & \textbf{$T_2$} & \textbf{Reference} \\
\midrule
Standard confocal (NA 0.9) & $\sim 5 \times 10^4$ & $\sim 1~\mu$s & \cite{Barry2020} \\  % Barry2020 Table 1 (confocal peak count rate)
Solid immersion lens (SIL) & $\sim 3.5 \times 10^5$ & preserved & \cite{Hadden2010} \\  % Hadden2010 Fig. 3 (SIL-enhanced count rate)
Bullseye grating & $(2.7 \pm 0.1) \times 10^6$ & 1.7 ms & \cite{Li2015} \\  % Li2015 Fig. 4b (peak count rate) and Fig. 3 (T2)
Hybrid nanoantenna & $\gtrsim 10^7$ (projected) & --- & \cite{Barbiero2025} \\  % Barbiero2025 Table 1, projected upper bound from 82% collection efficiency
\bottomrule
\end{tabular}
\end{table}

\noindent
\textbf{Bullseye grating} \cite{Li2015}: A circular bullseye pattern etched into
the diamond surface achieves $R \approx 2.7 \times 10^6$~cps from a single NV
center with $T_2 = 1.7$~ms preserved, representing a $\sim 50\times$
improvement over standard confocal.

\textbf{Hybrid nanoantenna} \cite{Barbiero2025}: Ultra-precise positioning of
nanodiamonds in hybrid metal-dielectric bullseye nanoantennas achieves a
measured collection efficiency of 82\% into NA~$= 0.5$, approaching unity for
NA~$> 0.8$.  At near-unity collection from the NV's intrinsic $\sim 13$~ns
excited-state lifetime (radiative rate $\sim 7.7 \times 10^7$~s$^{-1}$), the
theoretical upper bound on detected photon rate approaches $\sim 10^7$~cps
after accounting for quantum efficiency and branching ratios.


\subsection{ODMR Contrast $C_1$: $m_s = 0$ vs.\ $m_s = \pm 1$}

The standard fluorescence contrast between the bright ($m_s = 0$) and dark
($m_s = \pm 1$) states arises from spin-selective intersystem crossing (ISC)
through the singlet states \cite{Doherty2013}.

\begin{itemize}
  \item \textbf{CW ODMR:} $C_1 \approx 0.20$--$0.30$ (single NV, confocal)
    \cite{Doherty2013, Barry2020}. % Doherty2013 Fig. 6 / Barry2020 Sec. II.A p.4
  \item \textbf{Pulsed ODMR:} $C_1$ up to $\approx 0.30$ (theoretical limit
    for single NV) \cite{Barry2020}. % Barry2020 Eq. (6) and Sec. II.A
  \item \textbf{Laser threshold magnetometry:} Enhanced effective contrast via
    stimulated emission at high pump power \cite{Doherty2013}. % Doherty2013 Sec. 3.4
\end{itemize}

\noindent
Throughout this document we adopt $C_1 = 0.30$ as the state-of-the-art value. % Barry2020 Eq. (6); Doherty2013 Fig. 6


\subsection{Differential Contrast $C_2$: $m_s = -1$ vs.\ $m_s = +1$}
\label{sec:C2}

This is the critical parameter for \emph{direct} three-level readout.

\textbf{Physical origin.}  The ISC from the $^3E$ excited state to the singlet
$^1A_1$ is mediated by spin-orbit coupling.  At the level of the $C_{3v}$
symmetry of the NV center, the ISC rates from the excited-state orbital
branches depend on the spin projection $|m_s|$ but \emph{not} on the sign of
$m_s$ \cite{Goldman2015a, Batalov2009}.  The states $|A_1\rangle$ and
$|E_{1,2}\rangle$ in the excited-state fine structure have enhanced ISC rates
($\Gamma_{A_1}/2\pi \approx 16$~MHz), while the $|A_2\rangle$, $|E_x\rangle$,
$|E_y\rangle$ branches have suppressed ISC.  Crucially, $m_s = +1$ and
$m_s = -1$ enter the same orbital branches with equal weight; the ISC is
symmetric under $m_s \to -m_s$ \cite{Goldman2015a}.

\textbf{At $B = 0$:} $C_2 = 0$ exactly, by time-reversal symmetry.  The
$m_s = \pm 1$ states are degenerate and indistinguishable in fluorescence.

\textbf{At $B \neq 0$:} The Zeeman effect splits $m_s = \pm 1$ in frequency
space (observable as two ODMR dips), but the \emph{fluorescence intensity}
from each state remains identical to within experimental precision
\cite{Doherty2013}.  The ISC rate depends on $|m_s|$, not $m_s$.

\textbf{Possible mechanisms for nonzero $C_2$:}
\begin{enumerate}
  \item \textbf{Transverse strain/electric field:} Mixes the excited-state
    orbital branches, creating asymmetric ISC.  The effect is
    sample-dependent and typically small ($C_2/C_1 < 0.01$)
    \cite{Doherty2013}.
  \item \textbf{Off-axis magnetic field:} Mixes $m_s = 0$ and $m_s = \pm 1$
    in the excited state, creating slight fluorescence asymmetry.  At the
    20--50~G operating range along the NV axis, this is negligible.
  \item \textbf{Nuclear spin polarization:} At high $B$-fields near the
    excited-state level anti-crossing (ESLAC, $\sim 512$~G), nuclear-spin-dependent
    ISC creates measurable asymmetry.  We operate far from ESLAC.
\end{enumerate}

\textbf{Conclusion:} For all practical purposes at $B = 20$--$50$~G,
\begin{equation}
  \boxed{C_2 \approx 0.}
  \label{eq:C2zero}
\end{equation}
\emph{Direct fluorescence cannot distinguish $m_s = -1$ from $m_s = +1$.}
This is a fundamental constraint on qutrit readout architecture, not a
technology limitation.


\subsection{Single-Shot Readout Fidelity}

\begin{itemize}
  \item \textbf{All-optical (cryogenic resonant excitation):} $F > 95\%$
    single-shot, up to 99.9\% with spin-to-charge conversion (SCC)
    \cite{Zhang2021}.  % Zhang2021 Table 1 (cryogenic SCC fidelity) and abstract
    The SCC method maps the spin state to the charge state
    (NV$^-$ vs.\ NV$^0$), which can be read with near-perfect fidelity
    ($F_{\text{charge}} = 99.96\%$). % Zhang2021 Table 1, row: charge readout fidelity
  \item \textbf{Room-temperature (repeated readout):} $F \sim 0.01$ per
    single readout attempt (photon shot noise limited).  Multiple
    measurements average to improve fidelity.
  \item \textbf{Spin-to-charge conversion at RT:} $F \approx 98.1\%$ with
    inefficient optics; $> 99\%$ projected with photonic structures
    \cite{Zhang2021}. % Zhang2021 Table 1 (room-temperature SCC fidelity = 98.1%) and Sec. IV
\end{itemize}

\textbf{Key point for ensemble readout:}  An ensemble of $N$ NV centers
provides $\sqrt{N}$ improvement in signal-to-noise ratio (SNR) per readout
window, equivalent to averaging $N$ independent single-shot measurements.


\subsection{Coherence Times in Isotopically Purified $^{12}$C Diamond}

\begin{table}[h]
\centering
\caption{State-of-the-art electron spin coherence in isotopically purified diamond.}
\label{tab:coherence}
\begin{tabular}{lcccl}
\toprule
\textbf{Parameter} & \textbf{Value} & \textbf{$T$ (K)} & \textbf{$^{13}$C} & \textbf{Reference} \\
\midrule
$T_2^*$ & $\sim 90~\mu$s & 300 & 0.3\% & \cite{Maurer2012} \\   % Maurer2012 p.1284, Table 1
$T_2$ (Hahn echo) & 1.8 ms & 300 & 0.3\% & \cite{Balasubramanian2009} \\  % Balasubramanian2009 p.384, Fig. 2c
$T_2$ (DD, 77 K) & $\sim 0.6$ s & 77 & $0.01\%$ & \cite{BarGill2013} \\  % BarGill2013 Fig. 3, Table 1
$T_2$ (DD, RT) & approaches $2T_1$ & 300 & $0.01\%$ & \cite{Cambria2025} \\  % Cambria2025 Fig. 2
$T_1$ (cryo) & $> 1$ s & $< 50$ & --- & \cite{Jarmola2012} \\  % Jarmola2012 Fig. 3 (low-T plateau)
\bottomrule
\end{tabular}
\end{table}


%% ==================================================================
\section{Direct Three-Level Readout: Ensemble Size Derivation}
\label{sec:direct}
%% ==================================================================

\subsection{Setup}

Consider $N$ NV centers comprising one logical qutrit, in the state
\begin{equation}
  \ket{\psi} = a \ket{m_s = 0} + b \ket{m_s = -1} + c \ket{m_s = +1},
  \quad |a|^2 + |b|^2 + |c|^2 = 1.
\end{equation}
In a readout window of duration $T_r$, each NV center emits photons at rate
$R$ (cps).  The fluorescence depends on the spin state:
\begin{align}
  \text{Rate}(m_s = 0) &= R, \\
  \text{Rate}(m_s = -1) &= R(1 - C_1), \\
  \text{Rate}(m_s = +1) &= R(1 - C_1)(1 - C_2),
\end{align}
where $C_1 \approx 0.30$ is the standard ODMR contrast % Barry2020 Sec. II.A; Doherty2013 Fig. 6
and $C_2$ is the
differential contrast between the two dark states.

\subsection{Signal-to-Noise Ratio}

The expected photon count from the ensemble in window $T_r$ is:
\begin{equation}
  \langle n \rangle = N R T_r \big[ |a|^2 + (1-C_1)|b|^2 + (1-C_1)(1-C_2)|c|^2 \big].
\end{equation}
Photon shot noise gives $\sigma_n = \sqrt{\langle n \rangle}$.

To distinguish $m_s = -1$ from $m_s = +1$, we need to resolve the difference
in photon rates:
\begin{equation}
  \Delta n_{-1,+1} = N R T_r (1-C_1) C_2 \big( |b|^2 - |c|^2 \big).
\end{equation}
The SNR for this discrimination is:
\begin{equation}
  \text{SNR}_{-1,+1} = \frac{|\Delta n_{-1,+1}|}{\sqrt{\langle n \rangle}}
  \approx C_2 (1 - C_1)^{1/2} |b^2 - c^2| \sqrt{N R T_r},
  \label{eq:snr_direct}
\end{equation}
where we approximated the noise by $\sqrt{N R T_r}$ (shot noise dominated).

\subsection{Minimum Ensemble Size}

For a target SNR threshold $\text{SNR}_{\text{th}}$ (related to fidelity
$F_{\text{th}}$ via the error function: $F_{\text{th}} = \frac{1}{2}\text{erfc}(-\text{SNR}_{\text{th}}/\sqrt{2})$), and taking $|b^2 - c^2| \sim 1$
(worst case is $|b^2 - c^2| \to 0$, best case is 1):
\begin{equation}
  \boxed{N_{\min}^{\text{direct}} = \frac{\text{SNR}_{\text{th}}^2}{C_2^2 (1-C_1) R T_r}.}
  \label{eq:Nmin_direct}
\end{equation}

For 99\% fidelity, $\text{SNR}_{\text{th}} \approx 2.33$ (from the inverse
error function).  For 99.9\%, $\text{SNR}_{\text{th}} \approx 3.09$.

\subsection{Numerical Estimates}

Given $C_2 \approx 0$ (\cref{sec:C2}), direct three-level readout is
\textbf{fundamentally impossible} via fluorescence alone.  To illustrate the
scaling even with optimistic nonzero $C_2$:

\begin{table}[h]
\centering
\caption{$N_{\min}^{\text{direct}}$ for direct three-level readout at $F = 99\%$
($\text{SNR}_{\text{th}} = 2.33$), $T_r = 300$~ns, various $C_2$ assumptions.}
\label{tab:Nmin_direct}
\begin{tabular}{lccc}
\toprule
$C_2$ (assumed) & Confocal ($5 \times 10^4$) & Bullseye ($2.7 \times 10^6$) & Nanoantenna ($10^7$) \\
\midrule
$10^{-2}$ & $5.2 \times 10^9$ & $9.5 \times 10^7$ & $2.6 \times 10^7$ \\
$10^{-3}$ & $5.2 \times 10^{11}$ & $9.5 \times 10^9$ & $2.6 \times 10^9$ \\
$10^{-4}$ & $5.2 \times 10^{13}$ & $9.5 \times 10^{11}$ & $2.6 \times 10^{11}$ \\
\bottomrule
\end{tabular}
\end{table}

\noindent
Even the most optimistic $C_2 = 0.01$ requires $N > 10^7$ NV centers per
qudit.  At realistic $C_2 \lesssim 10^{-3}$, the numbers are absurdly large.
\textbf{Direct three-level readout via fluorescence is not viable.}


%% ==================================================================
\section{Sequential Binary Readout Protocol}
\label{sec:sequential}
%% ==================================================================

\subsection{Protocol Description}

Instead of attempting to distinguish $m_s = -1$ from $m_s = +1$ directly,
we perform two sequential measurements, each using only the large $C_1$
contrast:

\textbf{Step 1:} Apply a resonant MW $\pi$-pulse on the $|0\rangle \leftrightarrow |-1\rangle$ transition.  This maps:
\begin{align}
  |0\rangle &\to |-1\rangle, \quad |-1\rangle \to |0\rangle, \quad |+1\rangle \to |+1\rangle.
\end{align}
Fluorescence readout now yields bright signal proportional to $|b|^2$ (the
population originally in $|-1\rangle$, now mapped to $|0\rangle$) and dark
signal from $|a|^2 + |c|^2$.

\textbf{Step 2 (fresh copy or re-preparation):} Apply a resonant MW $\pi$-pulse
on $|0\rangle \leftrightarrow |+1\rangle$.  This maps:
\begin{align}
  |0\rangle &\to |+1\rangle, \quad |+1\rangle \to |0\rangle, \quad |-1\rangle \to |-1\rangle.
\end{align}
Fluorescence readout yields bright signal proportional to $|c|^2$.

\textbf{Population extraction:}
\begin{align}
  P_{-1} &= |b|^2: \text{ from Step 1 fluorescence (bright fraction)}, \\
  P_{+1} &= |c|^2: \text{ from Step 2 fluorescence (bright fraction)}, \\
  P_0 &= |a|^2 = 1 - |b|^2 - |c|^2: \text{ by normalization}.
\end{align}

\subsection{SNR Analysis}

Each step is a standard binary discrimination between $m_s = 0$ (bright,
rate $R$) and $m_s = \pm 1$ (dark, rate $R(1-C_1)$), with contrast $C_1 = 0.30$.

For $N$ NV centers in readout window $T_r$:
\begin{align}
  \text{Signal per step} &= N R T_r C_1 \cdot p, \\
  \text{Noise} &= \sqrt{N R T_r}, \\
  \text{SNR per step} &= C_1 \sqrt{N R T_r} \cdot p,
\end{align}
where $p$ is the population being measured (e.g., $p = |b|^2$ for Step 1).
For the worst case $p \sim 1/3$ (uniform qutrit):
\begin{equation}
  \text{SNR}_{\text{step}} \approx \frac{C_1}{3} \sqrt{N R T_r}.
\end{equation}

The total readout fidelity for all three populations is bounded by the
worst single-step fidelity.  Setting $\text{SNR}_{\text{step}} \geq \text{SNR}_{\text{th}}$:
\begin{equation}
  \boxed{N_{\min}^{\text{seq}} = \frac{9 \, \text{SNR}_{\text{th}}^2}{C_1^2 \, R \, T_r}.}
  \label{eq:Nmin_seq}
\end{equation}

\subsection{Comparison: Sequential vs.\ Direct}

The ratio of minimum ensemble sizes is:
\begin{equation}
  \frac{N_{\min}^{\text{direct}}}{N_{\min}^{\text{seq}}}
  = \frac{C_1^2}{9 \, C_2^2 (1 - C_1)}
  = \frac{(0.30)^2}{9 \, C_2^2 \times 0.70}
  = \frac{0.0143}{C_2^2}.
\end{equation}
For $C_2 = 10^{-2}$: ratio $= 143$.  For $C_2 = 10^{-3}$: ratio $= 1.43 \times 10^4$.

\textbf{Sequential binary readout is always superior} by a factor
$\sim C_1^2/(9 C_2^2)$, which diverges as $C_2 \to 0$.

\subsection{Numerical Results for Sequential Protocol}

\begin{table}[h]
\centering
\caption{$N_{\min}^{\text{seq}}$ for sequential binary qutrit readout.
$F = 99\%$ ($\text{SNR}_{\text{th}} = 2.33$), $C_1 = 0.30$.}
\label{tab:Nmin_seq}
\begin{tabular}{lccc}
\toprule
$T_r$ & Confocal ($5 \times 10^4$) & Bullseye ($2.7 \times 10^6$) & Nanoantenna ($10^7$) \\
\midrule
300 ns & $3.2 \times 10^6$ & $6.0 \times 10^4$ & $1.6 \times 10^4$ \\
$1~\mu$s & $9.7 \times 10^5$ & $1.8 \times 10^4$ & $4.9 \times 10^3$ \\
$10~\mu$s & $9.7 \times 10^4$ & $1.8 \times 10^3$ & $490$ \\
$100~\mu$s & $9.7 \times 10^3$ & $180$ & $49$ \\
\bottomrule
\end{tabular}
\end{table}

\textbf{Key result:} With bullseye grating collection and $T_r = 10~\mu$s
readout window, only $N \sim 1800$ NV centers are needed per logical qutrit
at 99\% fidelity.  With nanoantenna collection, this drops to $\sim 500$.
These are physically realizable ensemble sizes.

\textbf{Overhead:} The sequential protocol requires \emph{two} readout
windows per qutrit measurement (vs.\ one for direct), but each window
uses the same apparatus.  The $2\times$ time overhead is negligible compared
to the $>100\times$ reduction in $N_{\min}$.

\textbf{Note on re-preparation:} Steps 1 and 2 require either (a) two
identical copies of the state (measurement is destructive), or (b) weak
measurement with post-selection, or (c) ensemble averaging over many
experimental runs.  In the ensemble architecture with $N$ NV centers per
qudit, each run provides one measurement; the two steps are performed on
separate runs with identical state preparation.


%% ==================================================================
\section{Phase Information and Full Qutrit Tomography}
\label{sec:tomography}
%% ==================================================================

\subsection{What Fluorescence Cannot Measure}

Fluorescence readout measures \emph{populations} $\{|a|^2, |b|^2, |c|^2\}$
--- the diagonal elements of the density matrix $\rho$ in the $\{|0\rangle,
|-1\rangle, |+1\rangle\}$ basis.  It cannot access the \emph{off-diagonal}
(coherence) elements:
\begin{equation}
  \rho = \begin{pmatrix}
    |a|^2 & a b^* & a c^* \\
    b a^* & |b|^2 & b c^* \\
    c b^* & c a^* & |c|^2
  \end{pmatrix}.
\end{equation}
The phases $\arg(a/b)$, $\arg(b/c)$ (two independent phases for a qutrit)
are encoded in the off-diagonal elements.

\subsection{Measurement Bases for Full State Tomography}

A qutrit density matrix has $d^2 - 1 = 8$ independent real parameters
(3 populations minus 1 normalization constraint = 2, plus 3 complex
off-diagonal elements = 6 real parameters, total = 8).

To reconstruct all 8 parameters, one needs measurements in multiple bases.
The standard approach uses the \textbf{Gell-Mann matrices} $\{\lambda_1,
\ldots, \lambda_8\}$ as the SU(3) generators:

\begin{itemize}
  \item \textbf{Populations} (2 independent): Measured directly via
    fluorescence in the computational basis.
  \item \textbf{Coherences} (6 real parameters from 3 complex off-diagonal
    elements): Each requires measurement in a \emph{rotated basis}.
    Specifically, for each pair $(i,j)$, one needs measurements of both
    $\text{Re}(\rho_{ij})$ and $\text{Im}(\rho_{ij})$.
\end{itemize}

\textbf{Minimum measurement settings:} A complete set of measurements
requires at least \textbf{4 mutually unbiased bases} (each with 3 outcomes)
in dimension $d=3$, giving $4 \times 2 = 8$ independent real parameters
\cite{Wootters1989}.  In practice, \textbf{$\geq 8$ distinct measurement
settings} are used, each involving:
\begin{enumerate}
  \item A unitary rotation $U_k$ (implemented via MW pulses) that maps the
    desired measurement basis to the computational basis.
  \item Standard fluorescence readout in the computational basis (sequential
    binary protocol from \cref{sec:sequential}).
\end{enumerate}

\subsection{Scaling of Measurement Overhead}

For a single qutrit with sequential binary readout:
\begin{itemize}
  \item \textbf{Population measurement:} 2 readout windows $\times$ $M$
    repetitions for statistics.
  \item \textbf{Each coherence measurement:} 2 readout windows $\times$ $M$
    repetitions.
  \item \textbf{Total for full tomography:} $\geq 8$ measurement settings
    $\times$ 2 readout windows $\times$ $M$ repetitions = $16M$ readout
    windows.
\end{itemize}

For $n$ qudits ($d=3$), the density matrix has $d^{2n} - 1 = 3^{2n} - 1$
parameters.  Full tomography requires $O(d^{2n})$ measurement settings ---
exponential in $n$.  This is the standard quantum tomography scaling,
not specific to NV ensembles.

\textbf{Practical implication:} Full qutrit state tomography is feasible for
small qudit registers ($n \leq 3$--$4$) but becomes prohibitive for large
systems.  For quantum computing applications, one typically uses
process tomography or randomized benchmarking rather than full state
tomography.


%% ==================================================================
\section{Ensemble Size vs.\ Qudit Count Trade-Off}
\label{sec:tradeoff}
%% ==================================================================

\subsection{Architecture Parameters}

Let:
\begin{itemize}
  \item $M$ = total NV center budget (fixed by diamond sample)
  \item $k$ = NV centers per logical qudit (ensemble size)
  \item $Q = M/k$ = number of logical qudits
  \item $d$ = computational base (ternary: $d=3$)
  \item $F(k)$ = single-qudit readout fidelity as function of $k$
\end{itemize}

\subsection{Readout Fidelity as Function of $k$}

From \cref{eq:Nmin_seq}, the SNR per readout step for ensemble size $k$ is:
\begin{equation}
  \text{SNR}(k) = \frac{C_1}{3} \sqrt{k \, R \, T_r}.
\end{equation}
The single-qudit readout fidelity (for the worst-case population $p = 1/3$):
\begin{equation}
  F(k) = \frac{1}{2}\operatorname{erfc}\!\left(
    -\frac{C_1}{3\sqrt{2}} \sqrt{k \, R \, T_r}
  \right).
  \label{eq:Fk}
\end{equation}

\subsection{Computational Utility Metric}

The accessible Hilbert space dimension is $d^Q = d^{M/k}$, but errors
reduce the effective dimension.  We define a computational utility:
\begin{equation}
  \mathcal{U}(k) = Q \log_2 d \times F(k)^Q
  = \frac{M}{k} \log_2 d \times F(k)^{M/k}.
  \label{eq:utility}
\end{equation}
The first factor counts the number of useful qudits (in bits), while
$F(k)^Q$ is the probability that \emph{all} qudits are read correctly
(assuming independent errors).

Taking the logarithm for optimization:
\begin{equation}
  \ln \mathcal{U} = \ln\!\left(\frac{M \log_2 d}{k}\right)
  + \frac{M}{k} \ln F(k).
\end{equation}
Since $F(k) \to 1$ for large $k$, $\ln F(k) \approx -(1-F(k))$, and
\begin{equation}
  \ln \mathcal{U} \approx \ln\!\left(\frac{M \log_2 d}{k}\right)
  - \frac{M}{k} (1 - F(k)).
\end{equation}

\subsection{Optimal $k^*$}

The optimal $k^*$ balances two competing effects:
\begin{itemize}
  \item \textbf{Increasing $k$:} Improves $F(k)$ (better readout), but
    reduces $Q = M/k$ (fewer qudits).
  \item \textbf{Decreasing $k$:} More qudits, but each has worse fidelity.
\end{itemize}

The dominant term for large $M$ is $\frac{M}{k}(1 - F(k))$.  Substituting
$1 - F(k) \approx \frac{1}{\sqrt{2\pi}} \frac{3}{C_1 \sqrt{k R T_r}} e^{-C_1^2 k R T_r / 18}$
(Gaussian tail approximation), the exponent $\frac{M}{k} e^{-\alpha k}$
(with $\alpha = C_1^2 R T_r / 18$) vanishes for $k \gg 1/\alpha$ and
diverges for $k \ll 1/\alpha$.  Thus:
\begin{equation}
  \boxed{k^* \sim \frac{1}{\alpha} \ln M
  = \frac{18}{C_1^2 R T_r} \ln M.}
  \label{eq:kstar}
\end{equation}

\subsection{Numerical Estimates of $k^*$}

\begin{table}[h]
\centering
\caption{Optimal ensemble size $k^*$ for $M = 10^6$ total NVs,
$T_r = 10~\mu$s, $d = 3$.}
\label{tab:kstar}
\begin{tabular}{lcccc}
\toprule
\textbf{Collection} & $R$ (cps) & $\alpha$ (s$^{-1}$) & $k^*$ & $Q = M/k^*$ \\
\midrule
Confocal & $5 \times 10^4$ & $5 \times 10^{-3}$ & $\sim 2800$ & $\sim 360$ \\
Bullseye & $2.7 \times 10^6$ & $0.27$ & $\sim 51$ & $\sim 2 \times 10^4$ \\
Nanoantenna & $10^7$ & $1.0$ & $\sim 14$ & $\sim 7 \times 10^4$ \\
\bottomrule
\end{tabular}
\end{table}

\textbf{Qualitative answer:} The optimal $k^*$ is:
\begin{itemize}
  \item \textbf{Confocal:} $k^* \sim 10^3$ (regime of $\sim$thousands)
  \item \textbf{Bullseye:} $k^* \sim 50$ (regime of $\sim$tens)
  \item \textbf{Nanoantenna:} $k^* \sim 10$--$15$ (regime of $\sim$ten)
\end{itemize}
With state-of-the-art photon collection, the optimal ensemble is
surprisingly small --- of order 10--50 NV centers per logical qudit.


%% ==================================================================
\section{Practical Upper Bound on Qudit Base $d$}
\label{sec:maxd}
%% ==================================================================

\subsection{Generalization to Base $d$}

For a qudit of dimension $d$, the sequential binary readout protocol
requires $d-1$ binary measurements (each a $\pi$-pulse mapping one level
to $|0\rangle$ followed by fluorescence readout).  The worst-case population
is $p \sim 1/d$, giving:
\begin{equation}
  \text{SNR}_{\text{step}} = \frac{C_1}{d} \sqrt{k \, R \, T_r}.
\end{equation}
For target fidelity $F \geq 0.99$ on all $d-1$ binary steps simultaneously
(Bonferroni correction: require each step at $F_{\text{step}} \geq 1 - 0.01/(d-1)$):
\begin{equation}
  k_{\min}(d) = \frac{d^2 \, \text{SNR}_{\text{th}}^2(d)}{C_1^2 \, R \, T_r},
  \label{eq:kmin_d}
\end{equation}
where $\text{SNR}_{\text{th}}(d)$ increases logarithmically with $d$ due to
the Bonferroni correction.

\subsection{Maximum $d$ at Fixed $k$}

Inverting \cref{eq:kmin_d}:
\begin{equation}
  d_{\max} \approx \frac{C_1 \sqrt{k \, R \, T_r}}{\text{SNR}_{\text{th}}},
\end{equation}
where $\text{SNR}_{\text{th}} \approx 2.33$ for $F = 99\%$ (ignoring the
weak logarithmic $d$-dependence of the Bonferroni correction).

\begin{table}[h]
\centering
\caption{Maximum practical base $d$ at $> 99\%$ fidelity, for $k = k^*$ from
\cref{tab:kstar} and $T_r = 10~\mu$s.}
\label{tab:dmax}
\begin{tabular}{lccc}
\toprule
\textbf{Collection} & $k^*$ & $\sqrt{k^* R T_r}$ & $d_{\max}$ \\
\midrule
Confocal & 2800 & 37 & $\sim 5$ \\
Bullseye & 51 & 37 & $\sim 5$ \\
Nanoantenna & 14 & 37 & $\sim 5$ \\
\bottomrule
\end{tabular}
\end{table}

\textbf{Result:} All three technology tiers converge on $d_{\max} \approx 5$
at the optimal $k^*$, because $k^*$ is chosen precisely to give
$\sqrt{k^* R T_r} \approx \text{const}$ (it scales as $\ln M$).

However, this is the $d_{\max}$ \emph{at the optimal $k^*$ balancing qudit
count and fidelity}.  If we are willing to sacrifice qudit count for higher
$d$, we can increase $k$ beyond $k^*$:

\begin{table}[h]
\centering
\caption{$d_{\max}$ at $> 99\%$ fidelity for larger ensemble sizes.
$T_r = 10~\mu$s.}
\label{tab:dmax_largek}
\begin{tabular}{lcccc}
\toprule
\textbf{Collection} & $k = 100$ & $k = 1000$ & $k = 10^4$ & $k = 10^5$ \\
\midrule
Confocal & 3 & 9 & 29 & 91 \\
Bullseye & 21 & 67 & 211 & 669 \\
Nanoantenna & 41 & 129 & 407 & 1287 \\
\bottomrule
\end{tabular}
\end{table}

\textbf{Practical conclusion:} Ternary ($d = 3$) is achievable with very
modest ensemble sizes ($k \sim 50$--$100$ with bullseye collection).  Higher
bases $d = 5$--$10$ become practical with $k \sim 10^3$.  Going beyond
$d \sim 10$ requires either large ensembles or nanoantenna-class collection.

\textbf{Important caveat:} These bounds assume that the NV spin system can
be prepared in a $d$-level superposition.  The NV ground state is natively
a spin-1 system ($d = 3$).  Reaching $d > 3$ requires encoding into additional
degrees of freedom (e.g., nuclear spin, multiple NV orientations), which
introduces its own overhead and crosstalk issues beyond the scope of this
analysis.


%% ==================================================================
\section{Summary and Recommendations}
\label{sec:summary}
%% ==================================================================

\begin{enumerate}
  \item \textbf{$C_2 \approx 0$:} The differential fluorescence contrast
    between $m_s = -1$ and $m_s = +1$ is negligibly small at operating
    $B$-fields of 20--50~G.  Direct three-level readout via fluorescence
    is not viable.

  \item \textbf{Sequential binary readout is the correct protocol.}  Two
    measurements using $\pi$-pulse state mapping, each exploiting
    $C_1 \approx 0.30$, extract all three populations.

  \item \textbf{$N_{\min}$ for $d = 3$ at $F = 99\%$:}
    \begin{itemize}
      \item Confocal: $N \sim 10^5$ ($T_r = 10~\mu$s)
      \item Bullseye: $N \sim 1800$ ($T_r = 10~\mu$s)
      \item Nanoantenna: $N \sim 490$ ($T_r = 10~\mu$s)
    \end{itemize}

  \item \textbf{Optimal ensemble size:} With bullseye collection,
    $k^* \sim 50$ NV centers per qudit maximizes computational utility
    for a fixed NV budget of $M \sim 10^6$.

  \item \textbf{Maximum practical $d$:} At the optimal $k^*$,
    $d_{\max} \approx 5$.  Ternary ($d = 3$) is comfortably achievable.

  \item \textbf{Phase readout:} Full qutrit tomography requires $\geq 8$
    measurement settings (rotated bases via MW pulses), each consisting
    of two sequential binary readouts.  Overhead: $16M$ readout windows
    for $M$ statistical repetitions.

  \item \textbf{Technology recommendation:} Bullseye grating is the
    minimum viable photon collection technology for practical ensemble
    qutrit readout, achieving $k^* \sim 50$ per qudit with $> 99\%$
    fidelity.
\end{enumerate}


%% ==================================================================
%% BIBLIOGRAPHY
%% ==================================================================

\begin{thebibliography}{99}

\bibitem{Barry2020}
J.~F. Barry, J.~M. Schloss, E.~Bauch, M.~J. Turner, C.~A. Hart,
L.~M. Pham, and R.~L. Walsworth,
``Sensitivity optimization for NV-diamond magnetometry,''
\emph{Rev.\ Mod.\ Phys.} \textbf{92}, 015004 (2020).

\bibitem{Hadden2010}
J.~P. Hadden \emph{et al.},
``Strongly enhanced photon collection from diamond defect centers under
micro-fabricated integrated solid immersion lenses,''
\emph{Appl.\ Phys.\ Lett.} \textbf{97}, 241901 (2010).
arXiv:1006.2093.

\bibitem{Li2015}
L.~Li, E.~H. Chen, J.~Zheng, S.~L. Mouradian, F.~Dolde, T.~Schr\"oder,
S.~Karaveli, M.~L. Markham, D.~J. Twitchen, and D.~Englund,
``Efficient photon collection from a nitrogen vacancy center in a circular
bullseye grating,''
\emph{Nano Lett.} \textbf{15}, 1493--1497 (2015).
arXiv:1409.3068.

\bibitem{Barbiero2025}
A.~Barbiero \emph{et al.},
``Approaching unity photon collection from NV centers via ultra-precise
positioning of nanodiamonds in hybrid nanoantennas,''
\emph{APL Quantum} \textbf{2}, 036106 (2025).
arXiv:2503.14761.

\bibitem{Doherty2013}
M.~W. Doherty, N.~B. Manson, P.~Delaney, F.~Jelezko, J.~Wrachtrup,
and L.~C.~L. Hollenberg,
``The nitrogen-vacancy colour centre in diamond,''
\emph{Phys.\ Rep.} \textbf{528}, 1--45 (2013).

\bibitem{Goldman2015a}
M.~L. Goldman \emph{et al.},
``State-selective intersystem crossing in nitrogen-vacancy centers,''
\emph{Phys.\ Rev.\ B} \textbf{91}, 165201 (2015).
arXiv:1412.4865.

\bibitem{Batalov2009}
A.~Batalov \emph{et al.},
``Low temperature studies of the excited-state structure of negatively charged
nitrogen-vacancy color centers in diamond,''
\emph{Phys.\ Rev.\ Lett.} \textbf{102}, 195506 (2009).

\bibitem{Zhang2021}
Q.~Zhang \emph{et al.},
``High-fidelity single-shot readout of single electron spin in diamond with
spin-to-charge conversion,''
\emph{Nat.\ Commun.} \textbf{12}, 1529 (2021).

\bibitem{Maurer2012}
P.~C. Maurer \emph{et al.},
``Room-temperature quantum bit memory exceeding one second,''
\emph{Science} \textbf{336}, 1283--1286 (2012).

\bibitem{Balasubramanian2009}
G.~Balasubramanian \emph{et al.},
``Ultralong spin coherence time in isotopically engineered diamond,''
\emph{Nat.\ Mater.} \textbf{8}, 383--387 (2009).

\bibitem{BarGill2013}
N.~Bar-Gill, L.~M. Pham, A.~Jarmola, D.~Budker, and R.~L. Walsworth,
``Solid-state electronic spin coherence time approaching one second,''
\emph{Nat.\ Commun.} \textbf{4}, 1743 (2013).

\bibitem{Cambria2025}
M.~C. Cambria, S. Chand, C.~M. Reiter, and S. Kolkowitz,
``Scalable parallel measurement of individual nitrogen-vacancy centers,''
\emph{Phys.\ Rev.\ X} \textbf{15}, 031015 (2025).
arXiv:2408.11715.

\bibitem{Jarmola2012}
A.~Jarmola, V.~M. Acosta, K.~Jensen, S.~Chemerisov, and D.~Budker,
``Temperature- and magnetic-field-dependent longitudinal spin relaxation in
nitrogen-vacancy ensembles in diamond,''
\emph{Phys.\ Rev.\ Lett.} \textbf{108}, 197601 (2012).

\bibitem{Wootters1989}
W.~K. Wootters and B.~D. Fields,
``Optimal state-determination by mutually unbiased measurements,''
\emph{Ann.\ Phys.\ (N.Y.)} \textbf{191}, 363--381 (1989).

\end{thebibliography}

\end{document}
